\documentclass[11pt,a4paper,fontset=fandol]{ctexart}

\usepackage[a4paper,top=2.25cm,bottom=2.45cm,left=2.25cm,right=2.25cm,headheight=18pt,headsep=0.55cm,footskip=24pt]{geometry}
\usepackage{amsmath,amssymb,amsthm}
\usepackage{fancyhdr}
\usepackage{lastpage}
\usepackage{enumitem}
\usepackage{titlesec}
\usepackage{xcolor}
\usepackage{hyperref}
\usepackage{bookmark}
\usepackage[most]{tcolorbox}
\usepackage{setspace}
\usepackage{array}

\hypersetup{
  colorlinks=true,
  linkcolor=blue!50!black,
  urlcolor=blue!50!black
}

\setstretch{1.13}
\setlength{\parindent}{2em}
\setlength{\parskip}{0.25em}
\allowdisplaybreaks
\setlist[enumerate]{itemsep=0.45em, topsep=0.35em, leftmargin=2.4em}
\setlist[itemize]{itemsep=0.25em, topsep=0.25em, leftmargin=1.9em}

\definecolor{mainblue}{RGB}{36,83,140}
\definecolor{lightblue}{RGB}{241,246,252}
\definecolor{lightgray}{RGB}{246,246,246}
\definecolor{rulegray}{RGB}{110,118,128}

\titleformat{\section}
  {\Large\bfseries\color{mainblue}}
  {\thesection.}
  {0.45em}
  {}
  [\vspace{0.2em}\color{rulegray}\titlerule]
\titlespacing*{\section}{0pt}{1.1em}{0.7em}

\pagestyle{fancy}
\fancyhf{}
\fancyhead[L]{\small 数论方法中的组合构造周测 B 卷}
\fancyhead[C]{\small 组合专题：综合构造与不变量进阶}
\fancyhead[R]{\small 刘文博}
\fancyfoot[C]{\small 第 \thepage 页 / 共 \pageref{LastPage} 页}
\renewcommand{\headrulewidth}{0.55pt}
\renewcommand{\footrulewidth}{0.35pt}

\newtcolorbox{infobox}[1]{
  breakable,
  colback=lightblue,
  colframe=mainblue,
  boxrule=0.7pt,
  arc=1mm,
  title=\textbf{#1},
  fonttitle=\bfseries
}

\newenvironment{solution}{\par\noindent\textbf{参考解答：}}{\hfill$\square$\par}
\newcommand{\answerspace}{\vspace{2.3cm}}
\newcommand{\biganswerspace}{\vspace{3.2cm}}

\begin{document}

\thispagestyle{empty}
\begin{center}
{\fontsize{22}{28}\selectfont\bfseries 数论方法中的组合构造周测 B 卷\par}
\vspace{0.4cm}
{\large 适合小学高年级至初中低年级数学竞赛训练\par}
\end{center}

\begin{infobox}{考试说明}
\begin{itemize}
  \item 测试时间：75 分钟
  \item 满分：100 分
  \item B 卷更强调分类完整性和证明严密性，特别要注意：结论对“所有情况”是否都解释到了。
\end{itemize}
\end{infobox}

\section{试题}

\begin{enumerate}
  \item （12 分）从 $1$ 到 $9$ 中删去一个数，再把剩下的 $8$ 个数分成 $4$ 对，使每对两个数的和都是奇数。问删去哪几个数时可能？并给出理由。
  \answerspace

  \item （12 分）在数轴上，一个棋子从 $0$ 出发，每步只能向左或向右走 $1$ 格。它能否恰好走 $7$ 步到达 $4$？它能否恰好走 $8$ 步到达 $4$？
  \answerspace

  \item （12 分）有 $9$ 盏灯，开始时全灭。每次操作必须同时改变恰好两盏灯的状态（亮变灭，灭变亮）。问能否经过若干次操作后，恰好只有 $1$ 盏灯亮着？
  \answerspace

  \item （12 分）在一个 $5\times 5$ 的方格棋盘上，一只棋子从左下角出发，每次只能向上、下、左、右移动 $1$ 格。它能否恰好走 $8$ 步到达右上角？它能否恰好走 $9$ 步到达右上角？
  \biganswerspace

  \item （12 分）从 $1$ 到 $12$ 中选出 $4$ 个数，使它们的和能被 $4$ 整除，有多少种选法？
  \biganswerspace

  \item （12 分）把 $1$ 到 $12$ 分成 $4$ 组，每组 $3$ 个数。能否使每组的和都是 $3$ 的倍数？若能，请给出一种构造，并说明每组的余数结构只能有哪些可能。
  \biganswerspace

  \item （14 分）从 $1$ 到 $15$ 中最多能选多少个数，使得任意 $3$ 个被选出的数的和都不是 $3$ 的倍数？
  \biganswerspace

  \item （14 分）从 $1$ 到 $12$ 中任取 $6$ 个数，证明其中一定能选出 $3$ 个数，使这 $3$ 个数的和是 $3$ 的倍数。
  \biganswerspace
\end{enumerate}

\newpage
\section{详细解答}

\begin{enumerate}
  \item \begin{solution}
  要使每对的和都是奇数，就要求每一对都由“一个奇数和一个偶数”组成。

  在 $1$ 到 $9$ 中，奇数有 5 个，偶数有 4 个。

  如果删去一个偶数，那么剩下奇数 5 个、偶数 3 个，奇偶数个数不相等，就不可能分成 4 对且每对一奇一偶。

  如果删去一个奇数，那么剩下奇数 4 个、偶数 4 个，个数相等，就可以一奇一偶地配对。

  所以只有删去奇数时才可能，也就是删去
  \[
  1,3,5,7,9
  \]
  中的任意一个时都可能。
  \end{solution}

  \item \begin{solution}
  设向右走了 $r$ 步，向左走了 $l$ 步，则
  \[
  r+l=\text{总步数},\qquad r-l=4.
  \]

  对于 7 步，有
  \[
  r+l=7,\qquad r-l=4.
  \]
  相加得
  \[
  2r=11,
  \]
  这不可能，所以 7 步到达 4 不可能。

  对于 8 步，有
  \[
  r+l=8,\qquad r-l=4.
  \]
  相加得
  \[
  2r=12,
  \qquad r=6,
  \qquad l=2.
  \]
  这是可行的，所以 8 步到达 4 可以。
  \end{solution}

  \item \begin{solution}
  设当前亮灯数为 $L$。

  每次同时改变两盏灯状态时，亮灯数只可能：
  \begin{itemize}
    \item 增加 2；
    \item 减少 2；
    \item 不变。
  \end{itemize}

  因此亮灯数的奇偶性保持不变。

  初始时 $L=0$，是偶数。所以以后任何时候亮灯数都只能是偶数，不可能变成奇数。

  “恰好只有 1 盏灯亮着”对应 $L=1$，是奇数，因此不可能。
  \end{solution}

  \item \begin{solution}
  把左下角记作 $(0,0)$，右上角记作 $(4,4)$。

  从 $(0,0)$ 到 $(4,4)$ 的最短距离是
  \[
  4+4=8.
  \]
  所以恰好走 8 步是可以的，例如向右 4 步、向上 4 步。

  再看 9 步。

  每走一步，坐标和 $x+y$ 的奇偶性都会改变一次。起点 $(0,0)$ 的坐标和是 0，为偶数；终点 $(4,4)$ 的坐标和是 8，也为偶数。

  从偶数奇偶性回到偶数奇偶性，必须走偶数步，所以走 9 步不可能到达。

  结论：8 步可以，9 步不可以。
  \end{solution}

  \item \begin{solution}
  把 $1$ 到 $12$ 按模 $4$ 分类：
  \[
  0\text{ 类： }4,8,12;\quad 1\text{ 类： }1,5,9;\quad 2\text{ 类： }2,6,10;\quad 3\text{ 类： }3,7,11.
  \]

  选 4 个数且和是 4 的倍数，余数类型只可能是：
  \[
  (0,0,1,3),\ (0,0,2,2),\ (0,1,1,2),\ (0,2,3,3),\ (1,1,3,3),\ (1,2,2,3).
  \]

  分别计数：
  \[
  \binom32\binom31\binom31=27,
  \quad
  \binom32\binom32=9,
  \quad
  \binom31\binom32\binom31=27,
  \]
  \[
  \binom31\binom31\binom32=27,
  \quad
  \binom32\binom32=9,
  \quad
  \binom31\binom32\binom31=27.
  \]

  总数为
  \[
  27+9+27+27+9+27=126.
  \]

  所以共有
  \[
  126
  \]
  种选法。
  \end{solution}

  \item \begin{solution}
  可以。

  一种分法是
  \[
  (3,6,9),\ (1,4,7),\ (2,5,8),\ (10,11,12).
  \]
  它们的和分别为
  \[
  18,12,15,33,
  \]
  都是 3 的倍数。

  下面说明余数结构。

  三个数的和要是 3 的倍数，则这三个数除以 3 的余数和必须为 0（模 3）。因此一组 3 个数的余数结构只可能是
  \[
  (0,0,0),\ (1,1,1),\ (2,2,2),\ (0,1,2).
  \]

  所以任何可行分组中的每一组，都只能是这四种类型之一。
  \end{solution}

  \item \begin{solution}
  把被选数按模 3 分成 0 类、1 类、2 类。

  如果某一类里有 3 个或以上的数，那么从这类中任选 3 个，它们的余数和就是
  \[
  0+0+0,\quad 1+1+1,\quad \text{或}\quad 2+2+2,
  \]
  都是 3 的倍数，这是不允许的。

  所以每一类最多只能选 2 个。

  另外，如果三类中都至少有 1 个数，那么从三类中各取 1 个，余数和为
  \[
  0+1+2=3,
  \]
  也是 3 的倍数，这也不允许。

  因此三类中至少有一类必须为空。于是所有被选数只能落在两类中，而且每类最多 2 个，所以最多只能选
  \[
  2+2=4
  \]
  个。

  这个上界可以达到，例如选
  \[
  \{1,4,2,5\}.
  \]
  任取 3 个时，余数和只可能是 $1$ 或 $2$（模 3），都不会是 0（模 3）。

  所以最多能选 4 个数。
  \end{solution}

  \item \begin{solution}
  把这 6 个数按模 3 的余数分成 0 类、1 类、2 类。

  若某一类至少有 3 个数，则从这一类中任选 3 个，它们的和就是 3 的倍数。

  若每一类都至多 2 个数，由于总共有 6 个数，因此三类个数只能恰好是 $2,2,2$。这时从三类中各取 1 个，余数和就是
  \[
  0+1+2=3,
  \]
  仍是 3 的倍数。

  所以无论怎样取，必定能从中找到 3 个数，它们的和是 3 的倍数。
  \end{solution}
\end{enumerate}

\end{document}
